% !TeX root = ../main.tex
\section{Why look-ahead does not help here}
\label{sec:mechanism}

The outcome sections establish \emph{that} look-ahead stopped paying. This section asks what
it does to the training signal, using generation~3's per-candidate training logs --- every
branch point, every sampled candidate, its oracle score, and (at $\K{=}5$) its simulated
continuation. These are training-side analyses of generation~3 and are reported here to
interpret the outcome result; they are recorded in that experiment's own analysis
(\S\ref{sec:reproducibility-note}).

Look-ahead is supposed to improve the \emph{ranking} of candidates: siblings that look
promising but unwind badly should fall below siblings that compound. Four measurements say it
does not do this.

\subsection{It does not create more decisions to make}
Look-ahead changes what the oracle sees, not where the branches are. The number of branch
points is set by the trunk length and the minimum-context filter, both matched across arms.
Empirically the $\K{=}5$ arm builds no additional branch points --- so whatever it buys, it is
not more opportunities to choose.

\subsection{It rescales the reward signal without sharpening it}
\label{sec:spread}
Scoring a five-turn continuation rather than a single reply produces more dispersed scores:
within-group score spread grows by roughly $1.55\times$. That looks like a better signal --- a
wider gap between best and worst candidate should make the choice easier.

It is not, because \textbf{the margin and the standard deviation rise by the identical
factor}. The ratio margin$/$SD --- the quantity that actually determines whether the winner is
distinguishable from the pack --- is unchanged, and in every arm it sits at
${\approx}2.85$, the value expected from pure noise for the best of eight independent draws.
Look-ahead multiplies the units of the reward without changing how much of it is signal.

This also explains a side effect that could otherwise be mistaken for a benefit. The
preference filter $\tau = 0.1$ is an \emph{absolute} threshold, so a rescaled reward clears it
more often and the $\K{=}5$ arm yields more preference pairs. Rescaling the $\K{=}0$ arm's
margins by the observed SD ratio closes $44$--$87\%$ of that yield gap: most of the extra data
is a unit change, not extra information.

\subsection{More data from look-ahead does not help}
\label{sec:starvation}
That extra yield gives a direct test of the most obvious rival explanation --- that PTO's curve
flattens because it runs out of preference pairs, and that look-ahead ought to help by
supplying more. The generation-3 arms let us check it, because the $\K{=}5$ arm really did
train on more pairs at matched iterations: $568$ versus $475$ at iteration~6 ($1.2\times$) and
$689$ versus $400$ at iteration~7 ($1.7\times$).

It scored \emph{worse} at iteration~6 ($\dlt = +0.257$, the one Holm-significant cell in the
whole comparison) and tied at iteration~7. \textbf{Seventy percent more preference data bought
nothing measurable}, which retires data starvation as an explanation for either the flattening
curve or look-ahead's failure.

\subsection{At matched policy it adds no reward faithfulness}
\label{sec:faithfulness}
The stated purpose of look-ahead is to make the short training reward a more faithful proxy for
the full-session score the thesis actually evaluates. This is measurable: for training cuts at
each conversation depth, how well does the oracle's score on the cut rank-order the eventual
full-conversation score?

Pooled over all iterations, the $\K{=}5$ arm does look more faithful ($+0.052$). But that
comparison is confounded, because by iteration $t$ the two arms have different policies and are
generating different conversations. Controlling for it --- comparing the two arms at
\emph{matched policy}, at training iteration~1, where both are still branching from the same
untrained base (base \QQ{} $3.000$ versus $3.003$, $p = 0.987$) --- the advantage vanishes:
look-ahead leads in $11$ of $19$ depth bins, weighted mean $-0.005$, Wilcoxon $p = 0.59$.
\textbf{Zero.} The pooled advantage was the policy difference, not the lever.

\subsection{Putting it together}
Look-ahead in this setting changes the scale of the reward, not its information content. It
adds no branch points, no discriminative power per branch point, and no proxy faithfulness once
the policy is controlled for; the extra preference pairs it yields are a threshold artifact and
do not improve the result. Against that, it costs several times more per iteration, because
every candidate at every branch point requires $\K$ further simulated turns before it can be
scored.

That is a coherent explanation for generations~2 and~3 --- a lever that does nothing should
measure as nothing --- but on its own it would predict a null in generation~1 too, and
generation~1 shows a solid effect. The next section takes up that gap.

\subsection{A note on provenance}
\label{sec:reproducibility-note}
The outcome contrasts in \S\ref{sec:gen1}--\S\ref{sec:gen3} are computed by this chapter's own
script from the three generations' score files. The four training-side measurements in this
section are results of generation~3's existing training-signal analysis, which reads
per-generation logs this chapter does not otherwise touch; they are cited here rather than
recomputed, and \texttt{NUMBERS.md} points at the artifacts. The preference-pair counts in
\S\ref{sec:starvation} were re-verified against those artifacts for this chapter.
